\begin{abstract}
\singlespacing

How scientific ideas are communicated can shape how they are evaluated, disseminated, and ultimately integrated into cumulative knowledge. This paper studies whether author-team gender composition is associated with differences in scientific writing in economics. Using 9,117 abstracts from top general-interest economics journals and a 16-item LLM-based rubric, we first benchmark the LLM readability measure against Hengel's established readability metrics. We then examine mean differences across writing dimensions and estimate regressions with extensive journal, year, article, and author controls. Female-majority teams write more readable, less jargon-heavy, and more direct abstracts, and they generally provide more evidentiary support. They also make weaker novelty claims. By contrast, differences in hedging, qualifiers, assertiveness, and emotional valence are small or inconsistent. 

Extending the analysis to peer review, we find that, unlike Hengel's readability measure, LLM readability shows no differential gender effect from the review process. In cumulative-exposure tests, the aggregate preemptive-adaptation pattern Hengel documents for readability is largely absent for LLM scores, though patterns resembling that mechanism appear in three specific dimensions: acknowledgement of limitations, modal verb strength, and novelty-claim strength. We read these dimension-specific patterns as suggestive rather than definitive, since evaluation noise and the limited granularity of the LLM rubric remain plausible contributors.

\end{abstract}

\vspace{0.5em}
\noindent \textbf{Keywords:} Gender differences; scientific writing; readability; natural language processing; team composition.

\noindent \textbf{JEL Classification:} J16; C55; C81.
